\section{}
    \begin{frame}[allowframebreaks]{Frequency Modulated M\"obius (FMM)}
        
        
        \cite{rueda-2021-a-novel-wave-decomposition-for-oscillatory-signals} propose a wave decomposition method for oscillatory signals based on Frequency Modulated M\"obius (FMM) waves.
        
        \begin{equation}
            W(t)
            =
            M
            +
            A
            \cos
            \Biggl[
                \beta
                +
                2
                \arctan
                \Biggl(
                \omega
                \tan
                \Bigl(
                \frac{t-\alpha}{2}
                \Bigr)
                \Biggr)
                \Biggr]
        \end{equation}
        
        \begin{itemize}
            \item $W(t)$ is the value of the FMM wave at time $t$.
            \item $M$ is the baseline parameter. Bound: $M \in \mathbb{R}$. It controls the vertical shift of the FMM signal without changing the waveform shape, phase, or temporal deformation.
            
            \item $t$ is time. Bound: $t \in [0, 2\pi)$.
            
            \item $A$ is the amplitude parameter. Bound: $A \geq 0$. It controls the vertical scale of the waveform.
            
            \item $\alpha$ is the location parameter. Bound: $\alpha \in [0, 2\pi)$. It controls the horizontal position of the waveform on the time axis.
            
            \item $\beta$ is the phase parameter. Bound: $\beta \in [0, 2\pi)$. It controls the phase shift of the waveform. In other words, it determins for example relative to $\alpha$, from where the maximum point of the wave starts
            
            \item $\omega$ is the shape parameter. Bound: $\omega \in (0, 1]$. It controls the sharpness, skewness, and temporal deformation of the waveform. Originally it is taken in $(0, 1]$ but in \cite{rueda-2021-a-novel-wave-decomposition-for-oscillatory-signals} when it is close to 0, it is more like a spike and when it is near 1 then it is more like a sinus wave.
        \end{itemize}
        
        \alert{Why FMM to represent an action potential oscillatory wave?}\\
        This research uses variational autoencoders to encode a trace to the latent space of the parameters of an FMM wave but why?
        \begin{itemize}
            \item Representing different sensory data modality in a unified representation (i.e. both GPS and LIDAR turn into an FMM, each representing an action potential for a stimuli)
        \end{itemize}
        
    \end{frame}